% TOP_PROBLEM: {"id":30007553,"problem_number":"LOCAL-30007553","title":"General equilibrium with limited higher-order beliefs","category_id":21,"category":{"id":21,"name":"economics","display_name":"Economics","description":"Open economic theory statements, published research agendas, and proposed scoped specifications; question provenance and historical priority are qualified per record.","slug":"economics","order_index":21,"created_at":"2026-10-08T00:00:00Z"},"status":"research_question","external_url":null,"legacy_ids":[],"published":false,"statement":"Identify and estimate how imperfect information and limited higher-order reasoning change the aggregate multiplier in the specified linear coordination economy, and distinguish these mechanisms using randomized information and belief elicitation.","economics":{"original_id":42,"field":"Economic theory, equilibrium and social choice","kind":"identification","formalization_basis":"adapted_research_question","source_status":"research_agenda","priority_status":"earliest_found","priority_note":"This is the earliest source in which the relevant agenda was verified during this audit; absolute priority has not been established.","earliest_verified_reference":{"authors":["George-Marios Angeletos","Chen Lian"],"title":"Dampening General Equilibrium: Incomplete Information and Bounded Rationality","year":2022,"url":"https://economics.mit.edu/sites/default/files/publications/handbook-GE-NBER.pdf","doi":"10.3386/w29776","locator":"§6.2, printed p.27; §7 Discussion, pp.28–31","evidence_summary":"The chapter calls for eliciting relevant subjective beliefs and empirically distinguishing specifications of limited reasoning and incomplete information."},"current_reference":{"authors":["George-Marios Angeletos","Chen Lian"],"title":"Dampening General Equilibrium: Incomplete Information and Bounded Rationality","year":2022,"url":"https://economics.mit.edu/sites/default/files/publications/handbook-GE-NBER.pdf","doi":"10.3386/w29776","locator":"§6.2, printed p.27; §7 Discussion, pp.28–31","evidence_summary":"The chapter calls for eliciting relevant subjective beliefs and empirically distinguishing specifications of limited reasoning and incomplete information."},"formalization_note":"The linear economy, treatment, multiplier estimand, and identification target are an adaptation; they are not a verbatim conjecture from the handbook chapter."},"statusQualification":"Published question or research agenda verified at the cited locator. The mathematical specification is adapted for this catalogue and includes additional assumptions; source publication does not establish first-ever priority or certify this exact model remains unsolved. The linear economy, treatment, multiplier estimand, and identification target are an adaptation; they are not a verbatim conjecture from the handbook chapter.","statusReviewedAt":"2026-10-08"}
% FIRST_POSTED: 2026-10-08
% ENTRY_KIND: statement_only
% !TeX program = lualatex
\documentclass[11pt]{article}
\usepackage{fontspec}
\usepackage[a4paper,margin=25mm]{geometry}
\usepackage{amsmath,amssymb,mathtools}
\usepackage{xurl}
\usepackage[hidelinks]{hyperref}
\newcommand{\sref}[2]{\hyperlink{source:#1:#2}{[S#2]}}
\setlength{\emergencystretch}{3em}
\begin{document}
% problemId:30007553
\section*{General equilibrium with limited higher-order beliefs}\label{30007553}
\subsection{Definitions and mathematical statement}
Let a population of agents \(i\in[0,1]\) choose scalar actions in a linear coordination economy,
\[
 a_i=\theta_i+\alpha E_i[\bar a],\qquad \bar a=\int_0^1a_i\,di,\quad0<\alpha<1.
\]
Here \(\theta_i\) is the payoff-relevant fundamental and \(E_i\) denotes the agent's subjective conditional expectation. Under full information and consistent common beliefs, a common additive shock to fundamentals has aggregate multiplier \(M_0=(1-\alpha)^{-1}\). Let \(Z\) be a randomized public-information treatment about such a shock, and let \(b_i^{(k)}\) denote elicited beliefs about other agents' beliefs at order \(k\). Define \(M_Z=\partial E[\bar a\mid do(\theta+\eta),Z]/\partial\eta\) at \(\eta=0\). The research target is to identify how \(M_Z/M_0\) varies with belief precision, reasoning depth, and their population distributions. Establish conditions under which a finite collection of action observations and incentivized belief elicitations separately identifies incomplete information and truncated reasoning, or derive sharp observational equivalence sets when it cannot. Estimate the multiplier and its uncertainty under randomized information and stable structural payoffs, and evaluate predictions on untreated economies. Existing attenuation results are maintained as background; the unresolved agenda concerns empirically disciplined specifications and portable quantitative magnitudes.

\par\textbf{Formulation and scope.} The linear economy, treatment, multiplier estimand, and identification target are an adaptation; they are not a verbatim conjecture from the handbook chapter.

\par\textbf{Open-question provenance.} An explicit question or research agenda was verified at the source locator. This does not by itself certify that every later version remains unresolved \sref{30007553}{1}.

\par\textbf{Historical priority.} This is the earliest source in which the relevant agenda was verified during this audit; absolute priority has not been established.

\subsection{Short English statement}
Identify and estimate how imperfect information and limited higher-order reasoning change the aggregate multiplier in the specified linear coordination economy, and distinguish these mechanisms using randomized information and belief elicitation.

\subsection{Sources}
\begin{itemize}\raggedright
\item[\textnormal{[S1]}]\hypertarget{source:30007553:1}{} George-Marios Angeletos, Chen Lian. 2022. Dampening General Equilibrium: Incomplete Information and Bounded Rationality. §6.2, printed p.27; §7 Discussion, pp.28–31.
\url{https://economics.mit.edu/sites/default/files/publications/handbook-GE-NBER.pdf}.
DOI: 10.3386/w29776.
{\small\textit{Earliest verified question/agenda reference; historical priority qualified below. The chapter calls for eliciting relevant subjective beliefs and empirically distinguishing specifications of limited reasoning and incomplete information.}}
\end{itemize}
\end{document}
